\documentclass[11pt]{article}
\usepackage{qfffcalc}

\QFFFMetadata
  {$\mathrm{SU}(3)+5F$}
  {\frac12}
  {$\mathrm{SU}(5)\times\mathrm{U}(1)_B\times\mathrm{U}(1)_I$}
  {$\mathrm{SU}(5)\times\mathrm{SU}(2)\times\mathrm{U}(1)_q$}
  {degrees $2$--$10$}
  {Finite--UV branching and charge map}
  {Complete}

\begin{document}
\MakeQFFFTitle

\section{Conventions}
\[
\mathrm{SU}(5)\times\mathrm{SU}(2)
\longrightarrow
\mathrm{SU}(5)\times\mathrm{U}(1)_x,
\qquad
[a_1,a_2,a_3,a_4;b]_{5,2}
=[a_1,a_2,a_3,a_4]_5\otimes[b]_2,
\]
\[
[b]_2\longrightarrow\sum_{r=0}^{b}1_{x=b-2r}.
\]
The raw abelian basis is $(x,q)$.

\section{Finite and ultraviolet plethystic logarithms}
{\small
\begin{xltabular}{\linewidth}{@{}>{\centering\arraybackslash}p{0.05\linewidth} p{0.39\linewidth} X@{}}
\toprule
$d$ & finite PL & UV PL\\
\midrule
2 & $[0,0,0,0]_{5}\QFFFplus [1,0,0,1]_{5}$
& $[0,0,0,0;0]_{5,2}\QFFFplus [0,0,0,0;2]_{5,2}\QFFFplus [1,0,0,1;0]_{5,2}$\\
3 & $\beta^{-1}[0,1,0,0]_{5}\QFFFplus \beta[0,0,1,0]_{5}$
& $q^{-1}[0,0,1,0;1]_{5,2}\QFFFplus q[0,1,0,0;1]_{5,2}$\\
4 & $\QFFFminus [0,0,0,0]_{5}\QFFFminus [1,0,0,1]_{5}$
& $\QFFFminus 2[0,0,0,0;0]_{5,2}\QFFFminus [1,0,0,1;0]_{5,2}$\\
5 & $\QFFFminus \beta^{-1}\bigl([0,0,1,1]\QFFFplus [0,1,0,0]\QFFFplus [2,0,0,0]\bigr)_{5}$
$\QFFFminus \beta\bigl([0,0,0,2]\QFFFplus [0,0,1,0]\QFFFplus [1,1,0,0]\bigr)_{5}$
& $\QFFFminus q^{-1}[0,0,0,2;1]_{5,2}\QFFFminus 2q^{-1}[0,0,1,0;1]_{5,2}\QFFFminus q^{-1}[1,1,0,0;1]_{5,2}$
$\QFFFminus q[0,0,1,1;1]_{5,2}\QFFFminus 2q[0,1,0,0;1]_{5,2}\QFFFminus q[2,0,0,0;1]_{5,2}$\\
6 & $\QFFFminus \beta^{-2}[0,0,0,1]_{5}\QFFFminus [0,0,0,0]_{5}\QFFFminus \beta^{2}[1,0,0,0]_{5}$
& $\QFFFminus q^{-2}[0,1,0,1;0]_{5,2}\QFFFminus q^{-2}[1,0,0,0;2]_{5,2}\QFFFminus [0,0,0,0;0]_{5,2}$
$\QFFFminus [0,0,0,0;2]_{5,2}\QFFFminus [0,1,1,0;0]_{5,2}\QFFFminus [0,1,1,0;2]_{5,2}$
$\QFFFminus [1,0,0,1;2]_{5,2}\QFFFminus q^2[0,0,0,1;2]_{5,2}\QFFFminus q^2[1,0,1,0;0]_{5,2}$\\
\bottomrule
\end{xltabular}
}

\section{Raw branching of the generators}
\subsection*{Degree 2}
\begin{align*}
[0,0,0,0;0]_{5,2}
&\longrightarrow [0,0,0,0]_{5;(0,0)},\\
[0,0,0,0;2]_{5,2}
&\longrightarrow [0,0,0,0]_{5;(2,0)}
\QFFFplus [0,0,0,0]_{5;(0,0)}
\QFFFplus [0,0,0,0]_{5;(-2,0)},\\
[1,0,0,1;0]_{5,2}
&\longrightarrow [1,0,0,1]_{5;(0,0)}.
\end{align*}

\subsection*{Degree 3}
\begin{align*}
q^{-1}[0,0,1,0;1]_{5,2}
&\longrightarrow [0,0,1,0]_{5;(1,-1)}
\QFFFplus [0,0,1,0]_{5;(-1,-1)},\\
q[0,1,0,0;1]_{5,2}
&\longrightarrow [0,1,0,0]_{5;(1,1)}
\QFFFplus [0,1,0,0]_{5;(-1,1)}.
\end{align*}
The ordered charge pair is $(x,q)$.

\section{Charge map}
\[
(2,0)_{x,q}\longmapsto(0,1)_{B,I},
\qquad
(1,1)_{x,q}\longmapsto(-3,0)_{B,I}.
\]
\begin{equation}
\begin{pmatrix}B\\ I\end{pmatrix}
=
\begin{pmatrix}
0&-3\\[1mm]
\frac12&-\frac12
\end{pmatrix}
\begin{pmatrix}x\\ q\end{pmatrix},
\qquad
B=-3q,
\qquad
I=\frac{x\QFFFminus q}{2}.
\end{equation}
The inverse map is
\[
q=-\frac{B}{3},
\qquad
x=2I-\frac{B}{3}.
\]

\section{Physical branching through degree six}
{\small\input{data/branching_physical_low.tex}}

\BeginQFFFAppendices

\clearpage
\begin{landscape}
\section{Complete native finite and ultraviolet plethystic logarithms}
{\footnotesize\input{data/branching_native_pl_table.tex}}
\end{landscape}

\clearpage
\begin{landscape}
\section{Complete raw branching through degree ten}
{\footnotesize\input{data/branching_raw_full.tex}}
\end{landscape}

\clearpage
\begin{landscape}
\section{Complete physical branching through degree ten}
{\footnotesize\input{data/branching_physical_full.tex}}
\end{landscape}

\clearpage
\begin{landscape}
\section{Complete finite-versus-UV comparison}
{\footnotesize
\renewcommand{\arraystretch}{1.06}%
\input{data/branching_comparison_sections.tex}}
\end{landscape}

\end{document}
